\section{局部逆与沿解移动的坐标}\label{sec:inverse-moving-coordinates}
解的 $C^1$、$C^r$ 依赖性已经由逐次逼近和变分方程证明。下面把这一分析工具与坐标变换联系起来：先构造局部逆，再沿已知解消去常数项和线性项。
局部逆的证明只使用第~\ref{sec:picard-contraction}~节的压缩映射定理，不依赖后面的直化定理。
\begin{theorem}[由压缩映射得到局部逆]\label{thm:local-inverse}
若 $F\in C^1$，且 $DF(x_0)$ 可逆，则 $F$ 在 $x_0$ 附近有 $C^1$ 局部逆。
\end{theorem}
\begin{proof}
平移定义点、值域原点，再左乘 $DF(x_0)^{-1}$，只需研究
$F(0)=0,DF(0)=I$ 的情形。写 $F(x)=x+R(x)$，则 $R(0)=0,DR(0)=0$。
由 $DR$ 连续，可取闭球 $B_r$，使其中 $\|DR\|\leq\theta<1$。
球的凸性与微积分中值估计给出
\[
 |R(x)-R(z)|\leq\theta|x-z|,
\]
特别地 $|R(x)|\leq\theta r$。

给定 $|y|<(1-\theta)r$，定义 $A_y(x)=y-R(x)$。
于是 $|A_y(x)|\leq|y|+\theta r\leq r$，且其压缩常数为 $\theta$。
因此存在唯一 $G(y)\in B_r$ 满足 $G(y)+R(G(y))=y$。
对 $y,z$ 的不动点方程相减，得到
\[
 |G(y)-G(z)|\leq|y-z|+\theta|G(y)-G(z)|,
\]
所以 $G$ 为 Lipschitz 连续，常数不超过 $(1-\theta)^{-1}$。

还要证明逆可微。固定 $y$，令 $x=G(y)$、$u=G(y+v)-G(y)$。
由刚得的界，$|u|\leq(1-\theta)^{-1}|v|$。
Taylor 公式给出 
\[
 v=F(x+u)-F(x)=DF(x)u+o(|u|).
\]
$DF(x)=I+DR(x)$ 可逆，其逆由收敛 Neumann 级数给出，故
\[
 u=DF(x)^{-1}v+o(|v|).
\]
于是 
\[
 DG(y)=DF(G(y))^{-1}.
\]
右端连续，所以 $G\in C^1$；而 $F$ 在小球内的单射性还可由
$|F(x)-F(z)|\geq(1-\theta)|x-z|$ 直接验证。
这给出了局部互逆的两个映射。
\end{proof}

\paragraph{沿解移动坐标与消去线性部分。}
设 $\varphi$ 是已知解，令 $x=\varphi(t)+\xi$，则新的右端在 $\xi=0$ 处为零。
令 $A(t)=D_xf(t,\varphi(t))$，写新方程为
\[
 \dot\xi=A(t)\xi+R(t,\xi),
\]
其中 $R(t,\xi)=o(|\xi|)$。
若 $X'=A(t)X$、$X(t_0)=I$，令 $\xi=X(t)\eta$，则
\[
 \dot\eta=X(t)^{-1}R(t,X(t)\eta).
\]
新方程在 $\eta=0$ 的常数项和线性项同时为零。
由近似解 $\eta(t)\equiv\eta_0$ 的积分残差
$o(|\eta_0|)$，结合前述误差估计，得到真解
\[
 \eta(t)=\eta_0+o(|\eta_0|).
\]
回到原坐标即得初值导数由 $X(t)$ 给出，与本章
命题~\ref{prop:matrix-variation} 相衔接。
